% ---------------------------------------------------------------------
% Resumen y palabras clave
% ---------------------------------------------------------------------

\section*{Resumen}

\noindent\textbf{Título original}: Diseño e Integración de un Dron para Navegación Autónoma Basada en Visión.

La dependencia de los drones respecto a los sistemas de posicionamiento satelital (GNSS) supone una vulnerabilidad crítica en su operatividad. Como alternativa, la navegación basada en visión, como la Geolocalización Visual Cruzada (CVGL), permite estimar la posición de la aeronave emparejando imágenes de a bordo con bases de datos georeferenciadas (satelitales, aéreas). Sin embargo, la integración práctica de estas soluciones en equipos ligeros supone un reto complejo debido a las restricciones de hardware, peso y energía.

Este Trabajo de Fin de Grado se centra en el diseño, ensamblaje e integración de un dron personalizado, concebido como banco de pruebas para aplicar un algoritmo de localización por visión. El proyecto abarca el montaje de los componentes físicos, la configuración del controlador de vuelo y su intercomunicación con una computadora de borde. Para comprobar la viabilidad del sistema, se llevarán a cabo experimentos de vuelo reales orientados a la captura de datos. Posteriormente, se evaluará el algoritmo de CVGL sobre las imágenes obtenidas para analizar su precisión y rendimiento computacional, explorando su ejecución en tiempo real a bordo si las capacidades del hardware lo permiten. El objetivo es validar la integración del hardware y el software, estudiando la viabilidad técnica del sistema para asistir en la navegación en escenarios sin señal GPS.

\textbf{Palabras clave}: Integración de Sistemas; Vehículos Aéreos No Tripulados; Computación de Borde; Navegación Basada en Visión; Geolocalización Visual Cruzada.

% ---------------------------------------------------------------------
\clearpage

\section*{\textit{Abstract}}

\begin{ParrafoIngles}

\noindent\textbf{Title}: Design and Integration of a Drone for Vision-Based Autonomous Navigation.

The reliance of drones on satellite positioning systems (GNSS) poses a critical vulnerability in their operability. As an alternative, vision-based navigation, such as Cross-View Geo-Localization (CVGL), allows for estimating the aircraft's position by matching onboard images with georeferenced databases (satellite, aerial). However, the practical integration of these solutions into lightweight equipment presents a complex challenge due to hardware, weight, and power constraints.

This Bachelor's Thesis focuses on the design, assembly, and integration of a custom drone, conceived as a testbed to apply a vision-based localization algorithm. The project encompasses the assembly of the physical components, the configuration of the flight controller, and its communication with an edge computer. To verify the system's viability, real-world flight experiments will be conducted aimed at data capture. Subsequently, the CVGL algorithm will be evaluated on the acquired images to analyze its accuracy and computational performance, exploring its real-time onboard execution if hardware capabilities permit. The objective is to validate the hardware and software integration, studying the technical feasibility of the system to assist navigation in GPS-denied environments.

\textbf{Keywords}: System Integration; Unmanned Aerial Vehicles; Edge Computing; Vision-Based Navigation; Cross-View Geo-localization.

\end{ParrafoIngles}

% ---------------------------------------------------------------------
\clearpage

\section*{Resum}

\begin{otherlanguage}{catalan}

\noindent\textbf{Títol}: Disseny i Integració d'un Dron per a Navegació Autònoma Basada en Visió.

La dependència dels drons respecte als sistemes de posicionament satel·lital (GNSS) suposa una vulnerabilitat crítica en la seua operativitat. Com a alternativa, la navegació basada en visió, com la Geolocalització Visual Creuada (CVGL), permet estimar la posició de l'aeronau aparellant imatges de a bord amb bases de dades georeferenciades (satel·litals, aèries). No obstant això, la integració pràctica d'aquestes solucions en equips lleugers suposa un repte complex a causa de les restriccions de maquinari, pes i energia.

Aquest Treball de Fi de Grau se centra en el disseny, muntatge i integració d'un dron personalitzat, concebut com a banc de proves per a aplicar un algorisme de localització per visió. El projecte comprén el muntatge dels components físics, la configuració del controlador de vol i la seua intercomunicació amb una computadora de vora. Per a comprovar la viabilitat del sistema, es duran a terme experiments de vol reals orientats a la captura de dades. Posteriorment, s'avaluarà l'algorisme de CVGL sobre les imatges obtingudes per a analitzar la seua precisió i rendiment computacional, explorant la seua execució en temps real a bord si les capacitats del maquinari ho permeten. L'objectiu és validar la integració del maquinari i el programari, estudiant la viabilitat tècnica del sistema per a assistir en la navegació en escenaris sense senyal GPS.

\textbf{Paraules clau}: Integració de Sistemes; Vehicles Aeris No Tripulats; Computació de Vora; Navegació Basada en Visió; Geolocalització Visual Creuada.

\end{otherlanguage}

\medskip
\rule{\textwidth}{0.5pt}

\section*{Títulos registrados}

\begin{tabularx}{\textwidth}{>{\bfseries}p{0.23\textwidth}X}
Castellano & Diseño e Integración de un Dron para Navegación Autónoma Basada en Visión.\\
Inglés & Design and Integration of a Drone for Vision-Based Autonomous Navigation.\\
Valenciano & Disseny i Integració d'un Dron per a Navegació Autònoma Basada en Visió.\\
\end{tabularx}

% ---------------------------------------------------------------------
